\documentclass[twocolumn,linenumbers]{aastex631}

\usepackage{amsmath}
\usepackage{booktabs}
\usepackage{graphicx}

\begin{document}

\title{Host-Mass Dependence of Planetary Near-Commensurability Varies with Dynamical Order}
\shorttitle{Near-Commensurability by Dynamical Order}
\shortauthors{Yarbrough}

\author{Leslie Yarbrough}
\affiliation{Independent Researcher, Antioch, TN, USA}
\email{leslie.yarbrough@gmail.com}

\begin{abstract}
Adjacent planets frequently have period ratios near integer
commensurabilities, but chance proximity is expected within finite,
dynamically constrained architectures. We analyze 326 systems containing 1150
planets from the NASA Exoplanet Archive using a system-conditioned null that
preserves host mass, planet count, planet masses, and period boundaries while
randomizing interior periods subject to an eight-mutual-Hill-radius separation.
An exact conditional sampler provides 2000 realizations for every geometrically
feasible system. A broad 14-ratio classifier gives G-mass hosts
($0.9\leq M_\star<1.1\,M_\odot$) a matched residual of $-5.46$ percentage
points (nominal 95\% bootstrap interval $[-10.81,-0.15]$). Restricting the
classifier to the first-order ratios 2:1, 3:2, 4:3, and 5:4 gives a G-host
residual of $+0.09$ points ($[-4.44,+4.69]$). Global quadratic host-mass tests
are not significant for either definition. The null is feasible for 300
systems and geometrically infeasible for 26. Infeasible systems are more
commensurable overall than feasible systems (65.7\% versus 56.5\% under the
broad definition; 32.9\% versus 22.0\% for first-order ratios). Among G-mass
hosts, the corresponding means are 69.3\% versus 47.5\% and 37.0\% versus
18.0\%. The apparent G-host deficit is therefore conditional on
commensurability order and on the domain in which nonresonant stable random
packing is possible. Compact systems outside that domain form a distinct,
highly commensurable population that stability-filtered null analyses must
report separately.
\end{abstract}

\keywords{Exoplanet systems (484) --- Orbital resonances (1181) ---
Planetary system formation (1257) --- Stellar masses (1614) ---
Astronomy data analysis (1858)}

\section{Introduction}

Adjacent period ratios encode the combined effects of planet formation,
migration, dissipation, and later dynamical evolution. Kepler multi-planet
systems exhibit a deficit of very closely spaced pairs and asymmetries near
some low-order commensurabilities \citep{fabrycky2014,weiss2018}. Planet sizes
are also correlated within individual systems
\citep{weiss2018,millholland2021}. Their separations are typically of order
12--22 mutual Hill radii \citep{fang_margot_2013,puwu2015}. Planetary
architectures are therefore structured, but the presence of structure does
not imply that every pair near an integer ratio is dynamically resonant.

Confirmation of mean-motion resonance normally requires libration of a
resonant angle. Such information is unavailable for most cataloged systems.
Moreover, planets confined to a finite radial interval can fall near selected
integer ratios by chance, especially after stability requirements restrict the
allowed spacing distribution. Population studies must distinguish nominal
period-ratio proximity from both confirmed resonance and the background
produced by stable packing.

Near-commensurability also evolves with time. Young close-in systems show a
substantially larger prevalence of first-order commensurabilities than mature
systems \citep{dai2024}, while Galactic kinematics indicate that Kepler
systems near second-order commensurabilities are younger than comparable
noncommensurable systems \citep{hamer2024}. These results motivate separating
dynamical orders rather than treating every selected integer ratio as an
equivalent tracer of formation history.

This distinction becomes important when many target ratios are used. A
classifier containing several higher-order commensurabilities can cover a
substantial portion of period-ratio space. Its resulting fraction may remain a
useful descriptive statistic when compared with a matched null, but it cannot
be interpreted directly as the fraction of resonant pairs. Results should also
be tested using the dynamically strongest first-order commensurabilities.

We examine whether near-commensurability relative to stable random packing
depends on host-star mass. The analysis extends an initial broad-ratio study in
three ways. First, it contrasts the complete 14-ratio classifier with a
first-order-only definition. Second, it replaces finite rejection sampling with
an exact conditional sampler that gives every feasible system equal simulation
weight. Third, it explicitly analyzes systems for which the adopted eight-Hill
null is geometrically impossible. This last population is not a random loss:
it includes compact resonant systems such as GJ~876 \citep{rivera2010} and
Kepler-80 \citep{macdonald2016} and is substantially more commensurable than
the systems described by the null. Resonant-chain systems including
TRAPPIST-1 \citep{luger2017} and Kepler-223 \citep{mills2016} provide further
examples of architectures for which nominal ratios and multi-body phase
relations must be distinguished from generic pairwise packing.

\section{Data and Methods}

\subsection{Canonical Catalog and Sample}

We use the NASA Exoplanet Archive \citep{christiansen2025}
\texttt{pscomppars} table (dataset DOI: 10.26133/NEA13) downloaded on
2026 April 27. The frozen catalog contains 1182 planets around 336 hosts and is
archived with the analysis code. The query required catalog multiplicity
$N_{\rm pl}\geq3$ and \texttt{pl\_controv\_flag=0}. We then require each
analyzed system to retain at least three planets with valid periods after
applying $0.5\leq P\leq10^4$ days, a finite positive host mass, and
$M_\star<1.6\,M_\odot$. These cuts leave 326 systems, 1150 planets, and 824
adjacent pairs.

The archive supplies composite stellar parameters on each planet row; these
can differ slightly within a host. To reproduce the frozen analysis, the
primary calculation uses the first finite host-mass value in catalog order.
A robustness calculation using the within-host median gives the same
qualitative distinction between the broad and first-order classifiers.

We label four host-mass intervals
\begin{align*}
{\rm M}:&\ M_\star<0.6\,M_\odot, &
{\rm K}:&\ 0.6\leq M_\star<0.9\,M_\odot,\\
{\rm G}:&\ 0.9\leq M_\star<1.1\,M_\odot, &
{\rm F}:&\ 1.1\leq M_\star<1.6\,M_\odot.
\end{align*}
These labels denote mass bins rather than independently verified spectral or
evolutionary classifications.

\subsection{Commensurability Definitions}

For adjacent planets ordered by period, let $R=P_2/P_1$. For a target set
$\mathcal{R}$, define
\begin{equation}
\delta=\min_{r\in\mathcal{R}}\left|\frac{R}{r}-1\right|.
\end{equation}
A pair is near-commensurable when $\delta\leq0.03$. This statistic measures
proximity to a nominal integer ratio; it does not demonstrate resonant-angle
libration or resonance depth.

The broad target set is
\begin{multline}
\mathcal{R}_{14}=\{3\!:\!2,2\!:\!1,4\!:\!3,5\!:\!3,5\!:\!4,7\!:\!5,
8\!:\!5,9\!:\!5,\\
7\!:\!3,5\!:\!2,8\!:\!3,3\!:\!1,7\!:\!2,4\!:\!1\}.
\end{multline}
It includes commensurabilities of several dynamical orders. Its nonoverlapping
$\pm3\%$ windows cover 55.5\% of the linear interval and 56.3\% of the
logarithmic interval over $1\leq R\leq4$.

The first-order set is
\begin{equation}
\mathcal{R}_1=\{2\!:\!1,3\!:\!2,4\!:\!3,5\!:\!4\},
\end{equation}
whose windows cover 12.2\% linearly and 17.3\% logarithmically over the same
domain. We report both definitions throughout.

\subsection{System-Conditioned Stability Null}

For each observed architecture, the null preserves host mass, planet count,
the planet-mass sequence, and the innermost and outermost periods. Interior
periods are randomized in log-period space subject to adjacent separations
greater than eight mutual Hill radii, following the empirical long-term
packing scale used for Kepler multi-planet systems \citep{puwu2015}. Planet
masses use the catalog best mass
when available. When no catalog mass is available, radius $R_p$ in Earth
radii is converted to mass in Earth masses using
\begin{equation}
M_p=\begin{cases}
R_p^{2.5}, & R_p\leq1.0,\\
R_p^{2.0}, & 1.0<R_p\leq1.5,\\
R_p^{1.5}, & 1.5<R_p\leq4.0,\\
R_p^{1.2}, & 4.0<R_p\leq8.0,\\
R_p, & R_p>8.0.
\end{cases}
\end{equation}
Planets lacking both quantities receive a fallback mass of $5\,M_\oplus$.
These estimates are used only to construct the Hill-filtered null; they do not
enter the observed period-ratio classification.

For adjacent semimajor axes $a_1$ and $a_2$, define $r_a=a_2/a_1$ and
\begin{equation}
h_i=8\left(\frac{m_i+m_{i+1}}{3M_\star}\right)^{1/3}.
\end{equation}
The pairwise Hill condition can be written
\begin{equation}
\frac{2(r_a-1)}{r_a+1}\geq h_i,
\end{equation}
which gives the exact minimum semimajor-axis ratio
\begin{equation}
r_{a,\min}=\frac{2+h_i}{2-h_i},
\end{equation}
and minimum log-period gap
\begin{equation}
d_i=\frac{3}{2}\ln r_{a,\min}.
\end{equation}

Ordered log-uniform interior periods induce a
$\mathrm{Dirichlet}(1,\ldots,1)$ distribution of adjacent gaps conditional on
their sum. After subtracting the required $d_i$, the remaining log-period
extent is sampled directly from the same uniform simplex. This produces the
target distribution of log-uniform architectures conditioned on the Hill cut
without rejection attrition. We generate 2000 realizations per feasible
system using NumPy's PCG64 generator with seed 20260711. Every included system
must have exactly 2000 accepted realizations; this condition is checked before
the per-system output is written.

If
\begin{equation}
\sum_i d_i>\ln(P_{\rm out}/P_{\rm in}),
\end{equation}
no architecture satisfying the adopted null exists within the observed
boundaries. We classify such a system as \emph{null-infeasible}. The null is
feasible for 300 systems and infeasible for 26.

\subsection{Matched Residuals and Statistical Tests}

For every feasible system $j$ and classifier $k$, we calculate
\begin{equation}
\Delta f_{j,k}=f_{j,k}^{\rm obs}-
\left\langle f_{j,k}^{\rm null}\right\rangle.
\end{equation}
Each architecture contributes once, irrespective of its multiplicity. Mean
residual intervals are obtained by bootstrapping systems 20,000 times within
each mass bin. Each replicate samples $N_b$ systems with replacement from the
$N_b$ systems in bin $b$ and records the unweighted mean residual. The quoted
interval is the 2.5th--97.5th percentile range. Resampling systems rather than
adjacent pairs preserves intra-system dependence and prevents high-multiplicity
architectures from dominating the uncertainty estimate.

We test continuous host-mass curvature using
\begin{equation}
\Delta f=\beta_0+\beta_1M_\star+\beta_2M_\star^2+\epsilon,
\end{equation}
and quote a one-tailed test of $\beta_2>0$. We also permute residuals among host
masses 20,000 times while preserving the host masses, bin sizes, and residual
distribution. The permutation $p$-value is the fraction of shuffled samples
whose fitted quadratic coefficient is at least as positive as the observed
coefficient, with a plus-one correction in numerator and denominator. These
tests evaluate a global U-shaped dependence; a single bin interval excluding
zero is not equivalent to a significant global trend.

The eight bin-level intervals (four mass bins under each of two classifiers)
are descriptive, nominal 95\% intervals. No multiple-comparison correction is
applied across them. In particular, the broad-classifier G interval only
narrowly excludes zero and should not be interpreted as family-wise
significant evidence after accounting for the eight comparisons.

\section{Results}

\subsection{Global Classification Fractions}

Among all 824 eligible adjacent pairs, 57.3\% satisfy the broad classifier and
22.9\% satisfy the first-order classifier. The similarity between the broad
fraction and the 56\% domain coverage underscores why the broad fraction must
be interpreted relative to a matched null rather than as an absolute resonance
occurrence rate.

The 300 null-feasible systems contain 1054 planets and 754 adjacent pairs;
56.5\% of those pairs satisfy the broad classifier and 22.0\% satisfy the
first-order definition. The 26 null-infeasible systems contain 96 planets and
70 pairs, of which 65.7\% and 32.9\%, respectively, satisfy the two
classifiers.

\subsection{Matched Residuals Depend on Dynamical Order}

Figure~\ref{fig:residuals} and Table~\ref{tab:residuals} show the central
matched comparison. Under $\mathcal{R}_{14}$, G-mass hosts have the most
negative mean residual, $-5.46$ percentage points, with a bootstrap interval
that narrowly excludes zero. M- and F-mass residuals are positive but
uncertain, while the K-mass residual is negative and uncertain.

Under $\mathcal{R}_1$, the G residual becomes $+0.09$ points and is
indistinguishable from zero. The largest first-order positive residual occurs
for F-mass hosts, $+5.40$ points, but its interval also includes zero. Thus the
G-host result is not a general deficit of the strongest low-order
commensurabilities. It is specific to the broader ratio set.

The eight intervals in Table~\ref{tab:residuals} are nominal and uncorrected
for multiple comparisons. The broad G interval's slight exclusion of zero is
therefore a marginal, classifier-specific feature rather than a
multiplicity-adjusted detection.

\begin{figure*}
\centering
\includegraphics[width=\textwidth]{../figures/fig1_residuals.pdf}
\caption{Mean system-level observed-minus-null residuals for null-feasible
systems. Left: the broad 14-ratio classifier. Right: first-order ratios only.
Error bars are system-bootstrap 95\% intervals. The broad-classifier G residual
is marginally below zero, whereas its first-order residual is effectively
zero.}
\label{fig:residuals}
\end{figure*}

\begin{table*}
\centering
\caption{Matched Residuals for Null-feasible Systems}
\label{tab:residuals}
\begin{tabular}{lrrrrrrr}
\toprule
& & \multicolumn{3}{c}{Fourteen-ratio classifier} &
\multicolumn{3}{c}{First-order classifier}\\
\cmidrule(lr){3-5}\cmidrule(lr){6-8}
Bin & $N$ & Observed & Null & $\Delta f$ (95\% CI) & Observed & Null & $\Delta f$ (95\% CI)\\
\midrule
M & 40 & 67.08 & 64.58 & $+2.50$ [$-6.17,+10.81$] & 23.54 & 20.96 & $+2.58$ [$-4.71,+10.09$]\\
K & 114 & 54.37 & 57.04 & $-2.67$ [$-8.26,+2.90$] & 18.79 & 18.24 & $+0.55$ [$-3.50,+4.70$]\\
G & 105 & 47.52 & 52.98 & $-5.46$ [$-10.81,-0.15$] & 18.02 & 17.92 & $+0.09$ [$-4.44,+4.69$]\\
F & 41 & 58.47 & 57.75 & $+0.72$ [$-7.57,+8.50$] & 24.74 & 19.35 & $+5.40$ [$-2.21,+13.35$]\\
\bottomrule
\end{tabular}
\tablecomments{Fractions and residuals are percentage points. Each system
contributes once.}
\end{table*}

For the broad classifier, the continuous quadratic coefficient is positive but
not significant ($p=0.163$); the residual-permutation result is $p=0.168$.
For first-order ratios the corresponding values are $p=0.168$ and $0.167$.
The data therefore do not establish a global U-shaped host-mass dependence.

\subsection{Null-Infeasible Systems Are More Commensurable}

The null-infeasible population is not a random cross-section of the sample.
Table~\ref{tab:feasibility} and Figure~\ref{fig:feasibility} compare observed
system-level fractions. In K-, G-, and F-mass bins, infeasible systems have
higher mean fractions under both classifiers. The contrast is largest among
G-mass hosts: the nine infeasible systems have a mean broad fraction of 69.3\%
and first-order fraction of 37.0\%, compared with 47.5\% and 18.0\% for the 105
feasible G-host systems.

\begin{table}
\centering
\caption{Observed Means by Null Feasibility}
\label{tab:feasibility}
\begin{tabular}{lrrrrrr}
\toprule
& \multicolumn{3}{c}{Feasible} & \multicolumn{3}{c}{Infeasible}\\
\cmidrule(lr){2-4}\cmidrule(lr){5-7}
Bin & $N$ & Full & First & $N$ & Full & First\\
\midrule
M & 40 & 67.1 & 23.5 & 5 & 53.3 & 23.3\\
K & 114 & 54.4 & 18.8 & 8 & 66.3 & 30.8\\
G & 105 & 47.5 & 18.0 & 9 & 69.3 & 37.0\\
F & 41 & 58.5 & 24.7 & 4 & 66.7 & 33.3\\
\bottomrule
\end{tabular}
\tablecomments{Values are mean system-level percentages.}
\end{table}

\begin{figure*}
\centering
\includegraphics[width=\textwidth]{../figures/fig2_feasibility.pdf}
\caption{Observed near-commensurability fractions for systems inside and
outside the geometric domain of the $8R_H$ null. Infeasible K-, G-, and F-mass
systems are more commensurable under both definitions. Error bars are
system-bootstrap 95\% intervals; solid bars indicate null-feasible systems,
and hatched bars indicate null-infeasible systems. Subgroup sizes appear in
Table~\ref{tab:feasibility}.}
\label{fig:feasibility}
\end{figure*}

The infeasible list includes known compact resonant architectures. GJ~876
hosts an extrasolar Laplace configuration \citep{rivera2010} and is
infeasible because its required minimum log-period extent under the catalog
masses is 4.862, larger than its available extent of 4.161. For Kepler-80 the
corresponding values are 3.966 and 2.697; its interconnected three-body
resonances have been confirmed dynamically \citep{macdonald2016}. Their
exclusion is therefore not a
Monte Carlo failure: no realization satisfying the adopted pairwise criterion
can be placed between their preserved boundaries.

The adjacent-ratio distributions of the two populations are shown in
Figure~\ref{fig:ratios}. The infeasible population contains only 70 adjacent
pairs, so individual features should not be overinterpreted, but its greater
concentration near several first-order ratios is consistent with the summary
fractions above.

\begin{figure*}
\centering
\includegraphics[width=\textwidth]{../figures/fig3_ratio_distribution.pdf}
\caption{Observed adjacent period-ratio distributions for null-feasible and
null-infeasible systems over $1\leq P_2/P_1\leq4$. Vertical lines mark the
first-order ratios 5:4, 4:3, 3:2, and 2:1. The solid blue histogram denotes
null-feasible pairs, and the dashed red histogram denotes null-infeasible
pairs.}
\label{fig:ratios}
\end{figure*}

\section{Discussion}

\subsection{What the G-Host Result Does and Does Not Establish}

Within the domain of the $8R_H$ null, G-mass systems contain fewer pairs near
the broad target set than their matched randomized architectures. The result
is marginal at the bin level and persists when host masses are summarized by
their within-system median. It is nevertheless not evidence that G-mass hosts
are generally deficient in mean-motion resonance. The first-order residual is
almost exactly zero, and the global host-mass curvature tests are not
significant.

The broad definition includes higher-order ratios and labels more than half of
the $1$--$4$ ratio domain. Its G-host deficit therefore describes the detailed
distribution of nominal integer proximity, not a measured shortage of
resonant locks. Determining which ratios drive the deficit and whether they
have a dynamical origin requires order-resolved modeling and resonant-angle
information.

\subsection{The Feasibility Boundary Is Astrophysical}

The 26 infeasible systems reveal a limitation of stability-filtered packing
nulls. A pairwise Hill threshold is an approximate nonresonant stability
criterion, not a necessary condition for every multi-planet system. Resonant
phase protection can stabilize compact configurations that violate such a
threshold. Consequently, conditioning on null feasibility can select against
architectures in which resonance is most dynamically important.

This interpretation is consistent with the broader observational evidence
that commensurability is more prevalent at young ages and is depleted as
systems evolve \citep{dai2024,hamer2024}. The feasibility boundary may isolate
a population whose compact resonant structure has survived, although the
present catalog-level analysis cannot determine whether resonance supplies
the required phase protection in every individual system.

The appropriate interpretation is conditional: the matched residuals apply to
systems for which a randomized eight-Hill architecture exists. The infeasible
systems must be reported separately rather than treated as missing simulations.
Their elevated commensurability is itself a result, but it does not by itself
prove resonance-supported stability. Catalog mass uncertainties, eccentricity,
inclination, and multi-body effects can also change the relationship between
pairwise Hill spacing and long-term survival.

This population suggests a direct next test. N-body integrations using mass
posteriors should compare the stability of the observed infeasible
architectures with noncommensurable perturbations that preserve their compact
spacing. Libration analysis can then determine whether the systems excluded by
the packing null are preferentially protected by resonance.

\subsection{Catalog and Detection Limitations}

The archive combines transit, radial-velocity, timing, and other detections
with heterogeneous completeness. Host-mass bins differ in target populations,
planet sizes, and multiplicities. A planet-level transit-only check preserves
the qualitative order dependence but weakens the broad G deficit, whose
interval then includes zero. The present result should therefore be treated as
an exploratory catalog-level finding rather than a definitive occurrence-rate
measurement.

Composite stellar masses can also differ between planet rows belonging to the
same host. Twenty-two systems in the frozen catalog cross a bin boundary under
the available row-level values. Using the host median does not restore a
first-order G deficit, but future work should use a homogeneous host-star table
and propagate stellar and planetary parameter uncertainties through both the
Hill feasibility calculation and the null ensemble.

\section{Conclusions}

We compared near-commensurability in 326 multi-planet systems with a
system-conditioned stability null and find:

\begin{enumerate}
\item The 14-ratio classifier covers approximately 56\% of the period-ratio
domain from 1 to 4; its global classified fraction must therefore be interpreted
relative to a matched null, not as a resonance occurrence rate.

\item Among 300 systems admitting an $8R_H$ randomized architecture, G-mass
hosts have a broad-classifier residual of $-5.46$ percentage points, with a
95\% bootstrap interval of $[-10.81,-0.15]$.

\item The G-host first-order residual is $+0.09$ points
($[-4.44,+4.69]$). The apparent deficit is therefore dependent on the adopted
ratio set and is not present for 2:1, 3:2, 4:3, and 5:4 alone.

\item Neither classifier yields a significant global quadratic dependence on
host mass.

\item Twenty-six systems are geometrically outside the domain of the adopted
null. They are more commensurable overall than the feasible sample, particularly
among G-mass hosts, and include known compact resonant systems.
\end{enumerate}

Stable random packing describes only part of the observed population. For the
systems it can describe, approximately solar-mass hosts show no deficit in
first-order near-commensurabilities. The systems it cannot describe are among
the most compact and commensurable architectures in the catalog. Treating the
feasibility boundary as a result rather than a simulation failure provides a
clearer route toward identifying when resonance is merely nearby in period
space and when it may be essential to planetary stability.

\section{Acknowledgments}

The author used ChatGPT \citep{openai2026} and Claude
\citep{anthropic2026} for code review, statistical verification, and drafting
assistance. The language models are not authors; the author reviewed all
outputs and accepts responsibility for the analysis, interpretation, citation
accuracy, and final text. This research has made use of the NASA Exoplanet
Archive, which is operated by the California Institute of Technology, under
contract with the National Aeronautics and Space Administration under the
Exoplanet Exploration Program.

\facility{Exoplanet Archive}

\software{NumPy \citep{harris2020}, SciPy \citep{virtanen2020},
Matplotlib \citep{hunter2007}}

\section{Data Availability}

The frozen 2026 April 27 catalog, canonical analysis script, exact sampler,
per-system results, manuscript source, and figure-generation code are archived
at Zenodo (doi:10.5281/zenodo.21326602). The canonical analysis uses seed
20260711 and 2000 draws per feasible system. This corrected version supersedes
the rejection-sampling outputs in version 4.0
(doi:10.5281/zenodo.21230853).

\bibliography{references}

\end{document}
